% ===== §5 Praxis: the strategy of no-strategy =====
\section{Praxis: The Strategy of No-Strategy}
\label{p09:sec:praxis}

\noindent
The epistemology ended on a result that dictates the form of the practice: that sustaining a relation is a sustained, fallible, first-personal attention, not the application of a metric. This section draws out what such a practice can and cannot be. Its conclusion is, at first, disquieting---that there is, strictly, no \emph{strategy} for sustainability---but the disquiet is the doorway to the answer, and the answer is the one the whole paper has been circling: a \emb{strategy of no-strategy}, which is \emph{wu wei} given its precise and non-mystical content.

\subsection{Why sublimation cannot be aimed at}
\label{p09:sec:no-aim}

Lay the foregoing in a line and an uncomfortable necessity appears. The mark of the good circle is positive holonomy---sublimation; and holonomy is accrued along a path walked in the real, not cashed out by a sign (\xref{p09:sec:semiotics} will make this the heart of the semiotics). But any intervention that takes ``producing sublimation'' as its direct aim is attempting to produce, directly, the very thing that can only arise as the by-product of a walked path---which is the definition of settling failure. \emb{Any strategy that takes the production of sublimation as its direct end will, precisely thereby, destroy it.} The more one schemes to ``manage the relation into sustainability,'' the more one flattens the spiral with the instrumental reason---the control, the extraction, the direct wringing-out of phase---that the geometric criterion identifies as vicious. A strategic love is not love; it is settling. To aim at the holonomy is, in the geometry, a category error: one cannot ``walk toward'' a holonomy, one can only walk the path, and the phase accrues of itself.

\begin{claim}[The strategy of no-strategy]
Sustainability has no strategy, only a disposition. The good circle cannot be produced, only not-obstructed. Every ``technique for making a relation sustainable'' is, structurally, a forced action, an extraction, and destroys what it would preserve by cashing out a sublimation never walked. The only self-consistent practice is a negative, custodial one: not to do something so as to produce sublimation, but to refrain from doing what destroys the curvature under which sublimation arises of itself---not to extract, not to hasten, not to settle, not to counterfeit real presence with a cheap sign. \emb{\emph{Wu wei} is not the absence of practice but the redirection of practice from ``producing sublimation'' to ``guarding the curvature under which sublimation occurs.''}
\end{claim}

\subsection{Three domains, three insights}
\label{p09:sec:three-domains}

The disposition shows itself, with a different lesson, in three domains; their parallel is part of the evidence that the structure is genuinely universal and not forced from a single case.

\emc{The Daoist domain} contributes the ground: that the natural dynamics trends toward the good (\xref{p09:sec:trending}), so that letting-be is not a passive default but a trust grounded in the system's own positive curvature. The wound heals; the body floats; the relation recovers.

\emc{The commercial domain} contributes a sharpening, and a caution. The highest selling is, very often, not to sell. To press a sale is to emit a self-undermining signal: the very act of pressing betrays that ``my interest and yours diverge, and so I must persuade you''---and the harder one presses, the more loudly one announces that the exchange may not, in truth, serve the buyer. The credibility of a relation of exchange cannot, in the end, rest on the polish of the pitch (which is cheap, outsourceable, machine-generable); it rests on the unforgeable fact that the thing is genuinely good, the relation genuinely mutual---so that the customer \emph{comes of themselves}, and a mutually beneficial circulation arises that needs no forcing. Incentive-compatibility produces an endogenous positive curvature: when the parties' interests align in the cycle, the cycle perpetuates itself with little forcing, and \emph{wu wei}---not interfering with a mutually beneficial cycle that would continue of itself---is the fitting disposition. (The maintenance cost reappears here as a marker of curvature; but it is held, as \xref{p09:sec:no-metric} demanded, to the status of a fallible marker, never a criterion.)

\emc{The domain of the rose} contributes the sharpest lesson of all, and lifts the doctrine from ``forced action is superfluous'' to ``forced action is harmful.'' One who, dreading that the rose will wilt, waters it without cease, changes its water, tends it with an anxious love---that one cannot alter the fact that the flower must fall, and worse: he rots the root, and brings about the withering he feared. Here forced action does not merely fail to help; it \emph{manufactures the very outcome it dreads}. The fear of an end, made into action, produces the end. This is \zh{为者败之，执者失之}---``the one who acts upon it ruins it, the one who grasps it loses it'' (\emph{Daodejing} ch.~64)---and it cuts at once through capital (whose dread of the end of growth drives the extraction that exhausts its base), through the intimate bond (whose dread of loss drives the control that rots the relation), and through the Dao itself. \emb{The deepest unsustainability is born of the fear of unsustainability.}

\subsection{Wu wei is not negligence: the diagnostic anchor}
\label{p09:sec:not-negligence}

But \emph{wu wei} must be sharply distinguished from negligence, or the doctrine becomes a high-toned alibi for abandonment---and the Daoist tradition has, historically, been put to just this use, sanctifying a ruler's inaction as cosmic wisdom. The paper must therefore furnish a criterion that divides the two; and the criterion is, exactly, the sign of the curvature.

\begin{claim}[Wu wei versus negligence]
\emph{Wu wei} is non-intervention in a positive-curvature system; negligence is non-intervention in a negative-curvature one. The same ``doing nothing'' guards a sublimation arising of itself in the first case, and tolerates an extraction underway in the second. Whether one ought, at this moment, to let be or to act is not settled by the feel of one's own disposition but by the diagnosis of the curvature: is this cycle, now, sublimating (returning to its root while the fibre accrues a positive phase), or is it being exhausted (a net extraction, one party giving without cease and nothing returning)? The legitimacy of \emph{wu wei} is local, and stands in need of ceaseless re-diagnosis; it is not a posture adopted once for all. The root of the rose may be left to its own regeneration; the partner being drained dry may not be left to ``wu wei'' her way to depletion---there, what is called for is a repairing action.
\end{claim}

\subsection{Repair is subtraction}
\label{p09:sec:repair}

What, then, is the repairing action that a negative-curvature system calls for? Not a new forced action laid on top of the old, but the \emph{removal} of the forced action that is doing the harm. This is the reconciliation of repair with \emph{wu wei}, and it is the section's final claim.

\begin{claim}[All true repair is subtraction]
Repair is not the addition of a new intervention but the subtraction of the one that sickens. The drowning body does not need to be taught a more elaborate stroke; it needs to be brought to cease its thrashing. The wound does not need a fiercer medicine; it needs the picking hand withdrawn. The intervention a negative-curvature system requires is, at bottom, the withdrawal of the anxious, extracting, force-closing hand---so that the system's own positive curvature may resume its work. \emb{The highest repair is itself a kind of \emph{wu wei}: not to do more, but to take away the hand that has all along been picking at the wound.} So \zh{扶}---to ``steady'' the drowning---does not mean to seize and haul him up (that is forced action, complicit in the thrashing); it means to hold him so that he ceases to struggle, and the water bears him up of itself. To steady is the least action that makes way for self-healing.
\end{claim}

The two dispositions are thus unified. \emph{Wu wei}, in a positive-curvature system, guards the sublimation arising of itself; repair, in a negative-curvature one, withdraws the hand that has bent the curvature negative---and thereby restores the conditions under which \emph{wu wei} is once more legitimate. Repair completes itself by giving way again to non-action. The practice of sustainability is, in sum, neither management nor abandonment but a vigilant, custodial attention: to keep reading the curvature, to stand back where it sublimates of itself, to withdraw the harming hand where it is being drained, and never to try to cash out, with a sign, the sublimation that can only be walked. It is neither a strategy nor a laissez-faire, but a poetic custody---like the reading of a poem that is never read out.
